\chapter{Conclusion and Outlook}
\label{ch:conclusion_outlook}

\section{Conclusion}

This work presented the Multi-Robot Calibration (MRC) framework, a
sensor-based system for determining the relative poses between robots
in a fleet using only onboard LiDAR sensors, without requiring
external infrastructure or manual measurement.
The system was implemented as a distributed ROS\,2 architecture
comprising approximately 3\,100 lines of C++, with a master--slave
communication pattern, two interchangeable registration backends
(KISS-Matcher and FRICP), optional ICP refinement, and GTSAM-based
pose graph optimization with star and mesh topologies.

The experimental evaluation, comprising over 700 individual calibration
runs across fleet sizes of 2, 3, 5, and 10 robots in 15 spatial
configurations, yielded the following key findings:

\begin{enumerate}
    \item \textbf{The optimal registration backend depends on the
          deployment scenario.}
          KISS-Matcher, a feature-based global registration method,
          consistently outperformed FRICP in the multi-robot distributed
          scenario ($N \geq 3$), where large inter-robot distances and
          diverse orientations cause FRICP to converge to local minima
          (rotation errors $> 120\degree$).
          In the two-robot proximity scenario, the result reversed:
          FRICP achieved sub-centimeter accuracy
          ($\bar{e}_t = \SI{5.4}{\milli\metre}$, $\bar{e}_r =
          0.03\degree$) by exploiting the favorable initial alignment
          and high point cloud overlap ($67\%$), while KISS-Matcher was
          limited to $\sim$\,\SI{12}{\centi\metre} by its feature
          quantization.

    \item \textbf{ICP refinement and pose graph optimization are not
          universally beneficial.}
          Averaged across all multi-robot experiments, both ICP
          refinement and pose graph optimization degraded accuracy.
          ICP shifted globally correct feature-based alignments toward
          locally dense regions ($+18\%$ RMS for KISS-Matcher), and
          pose graph optimization propagated errors from unreliable
          edges ($+16\%$ RMS for mesh topology).
          However, in favorable setups with high overlap and correct
          coarse alignment, ICP improved accuracy by up to $65\%$ and
          mesh PG by up to $80\%$.
          In the proximity scenario, ICP dramatically improved
          KISS-Matcher ($-86\%$) but degraded the already near-optimal
          FRICP result.

    \item \textbf{Sub-centimeter accuracy is achievable for the primary
          use case.}
          The two-robot proximity scenario---the most common real-world
          deployment for extrinsic calibration---achieved consistent
          sub-centimeter accuracy across all tested spatial arrangements
          (\SIrange{0.8}{2.0}{\metre} separation).

    \item \textbf{The modular software architecture enables systematic
          evaluation.}
          The factory pattern for registration backends, the runtime
          parameter reconfiguration via ROS\,2 services, and the
          automated evaluator node allowed all $2 \times 3 \times 2 = 12$
          parameter combinations to be tested within a single simulation
          session, enabling the comprehensive comparison presented in
          this work.
\end{enumerate}

\section{Outlook}

\subsection{Middleware Scalability}

The 10-robot experiments were limited by Zenoh middleware transport
saturation, which prevented systematic evaluation at larger fleet sizes.
Tuning Zenoh buffer sizes, using alternative DDS implementations, or
implementing a hierarchical communication architecture could enable
reliable operation with 10+ robots.
The partial 10-robot results showed that individual robots with
sufficient overlap to the master achieved sub-centimeter calibration,
indicating that the MRC algorithm itself scales but the communication
infrastructure requires adaptation.

\subsection{Pose Graph Improvements}

\paragraph{Efficient edge selection.}
The efficiency and accuracy of the pose graph can be improved by
selectively computing only pairwise registrations that are likely to
succeed.
Integration with UWB localization could provide rough positional
estimates for spatial filtering, reducing the computational burden of
a full-mesh calculation ($\binom{N}{2}$ edges) and avoiding the poor
results caused by low sensor overlap in high-offset scenarios.

\paragraph{Temporal constraints.}
Incorporating the time dimension by tracking robot trajectories would
increase edge density within the pose graph, providing more redundant
constraints and improving the reliability of the optimized result.

\paragraph{Multi-modal fusion.}
Additional sensing modalities such as UWB range measurements can
introduce further edges to the pose graph beyond LiDAR-based
registration, increasing connectivity and robustness.

\subsection{Registration Improvements}

UWB localization can also serve as a provider for initial
transformation estimates.
The evaluation showed that FRICP achieves superior accuracy when
the initial alignment is close to the true solution (proximity
scenario), but fails when starting from the identity transform
(distributed scenario).
Providing UWB-based initial guesses could extend FRICP's applicability
to the distributed scenario, potentially combining the sub-centimeter
precision of dense ICP with the robustness of having a good starting
point.

\subsection{Adaptive Configuration Selection}

The evaluation demonstrated that the optimal parameter combination
(backend, pose graph topology, ICP refinement) varies with fleet
size, spatial geometry, and point cloud overlap.
An adaptive strategy that automatically selects the registration
backend based on estimated inter-robot distance and overlap could
improve accuracy without manual tuning---for example, using FRICP
for nearby robot pairs and KISS-Matcher for distant ones.
